\section{Six geometries and their complete
census}\label{six-geometries-and-their-complete-census}

This application takes the odd branch of Corollary C.Odd. Retain the GAP
TomLib J1 permutation model, of order 175560. The construction
certificate constructs a marking \(a,b,c,d\) within that model with

\[
|ab|,|ac|,|ad|,|bc|,|bd|,|cd|=(5,2,2,3,2,5),
\] \[
(abc)^5=((dcb)^5a)^3=1.
\tag{J1}
\]

It also completes a Todd--Coxeter calculation for the literal
presentation with the four involution relations and (J1): the index over
\(\langle b,c,d\rangle\) is 1463. That subgroup is a quotient of the
finite Coxeter group {[}3,5{]} of order 120, so the presented group has
order at most \(1463·120=175560\). The found generating tuple in the
TomLib model supplies a quotient of that order; hence this is the source
presentation of J1. A valid permutation action alone would not prove the
upper bound: the completed exact coset enumeration is a charged
computational input.

Set \(w=(bcd)^5\). The certificate proves that w is a nonidentity
central involution of \(\langle b,c,d\rangle\). The six source-defined
markings and their exact invariants are:

{\small % Uncaptioned data table.
\begin{longtable}[]{@{}
  >{\raggedright\arraybackslash}p{(\linewidth - 6\tabcolsep) * \real{0.2500}}
  >{\raggedright\arraybackslash}p{(\linewidth - 6\tabcolsep) * \real{0.2500}}
  >{\raggedright\arraybackslash}p{(\linewidth - 6\tabcolsep) * \real{0.2500}}
  >{\raggedright\arraybackslash}p{(\linewidth - 6\tabcolsep) * \real{0.2500}}@{}}
\toprule\noalign{}
\begin{minipage}[b]{\linewidth}\raggedright
Geometry
\end{minipage} & \begin{minipage}[b]{\linewidth}\raggedright
Marking
\end{minipage} & \begin{minipage}[b]{\linewidth}\raggedright
Pair orders ab,ac,ad,bc,bd,cd in that marking
\end{minipage} & \begin{minipage}[b]{\linewidth}\raggedright
Maximal parabolic orders
\end{minipage} \\
\midrule\noalign{}
\endhead
\bottomrule\noalign{}
\endlastfoot
\(\Gamma_1\) & \(a,b,c,d\) & 5,2,2,3,2,5 & 120,20,20,60 \\
\(\Gamma_2\) & \(a,wb,c,d\) & 5,2,2,6,2,5 & 120,20,20,660 \\
\(\Gamma_3\) & \((a,b,c,c^d)\) & 5,2,2,3,3,5 & 120,20,60,60 \\
\(\Gamma_4\) & \((b^a,b,c,d)\) & 5,3,2,3,2,5 & 120,120,20,60 \\
\(\Gamma_5\) & \((b^a,b,c,c^d)\) & 5,3,3,3,3,5 & 120,120,60,60 \\
\(\Gamma_6\) & \(((wb)^a,wb,c,d)\) & 5,6,2,6,2,5 & 120,120,20,660 \\
\end{longtable}
}

Order-120 parabolics are \(C_{2}\times A_{5}\), order-20 ones are
\(D_{20}=C_{2}\times D_{10}\), order-60 ones A5 and order-660 ones
\(\operatorname{PSL}_{2}(11)\); D10 here has order ten. Structure
identifications use the retained exact GAP algorithms and named group
models.

\textbf{Corollary C.J1.} Each of these six marked coset systems is a
thin regular residually connected rank-four geometry with full
type-preserving group J1. They are pairwise nonisomorphic, even allowing
type permutations. These are {[}HL05{]} Theorems 3.2, 4.2, 5.2, 5.4, 6.2
and 7.2, including their generation lemmas.

\textbf{Proof.} Lemma C.Twist gives the central twist
\(\Gamma_1\to \Gamma_2\) and preserves its first parabolic. Every
endpoint operation has odd endpoint order five, so Corollary C.Odd
recovers the same group. The construction certificate verifies the two
actual subgroup equalities

\[
\langle b,c,d\rangle\cap\langle b,c,d\rangle^a=\langle c,d\rangle,
\qquad
\langle a,b,c\rangle\cap\langle a,b,c\rangle^d=\langle a,b\rangle.
\tag{J2}
\]

They have order ten. The second gives \(\Gamma_3\), the first
\(\Gamma_4\). In \(\Gamma_3\),
\(\langle b,c,c^d\rangle =\langle b,c,d\rangle\) because \(|cd|=5\), and
\(\langle c,c^d\rangle =\langle c,d\rangle\); hence the same first
equality gives \(\Gamma_5\). In \(\Gamma_2\) its first parabolic is
unchanged by the twist, so that equality gives \(\Gamma_6\). Thus after
the seed and twist cases, four IP proofs reduce to the two meets (J2).
The construction certificate additionally verifies all 256 ordered
marked-subset meets for each of the six markings, so the finite result
does not depend on a mistaken reuse of a subgroup identity.

For each marking it checks all sixteen equalities (C3), without
presuming that a C-group with a nonstring diagram is a geometry. Lemma
C.HellyCount and UC.Geometry now give all the geometric conclusions. The
six multisets of maximal parabolic orders in the table are distinct, so
no type-permuting isomorphism identifies two geometries. \(\square\)

The definition of \(\Gamma_3\) in {[}HL05{]} §5 duplicates an A5
parabolic, whereas its printed figure labelled \(\Gamma_3\) duplicates
\(C_{2}\times A_{5}\); the definition of \(\Gamma_4\) has the reverse
discrepancy. The source figure labels \(\Gamma_3\) and \(\Gamma_4\) are
therefore transposed relative to its definitions. The table follows the
definitions. This is a source correction, not a seventh construction.

The two source-mentioned failures also follow from the same certificate.
The markings \((a,wb,c,c^d)\) and (a,b,c,wd) still generate J1 and
satisfy typed Helly, but each has eight failing ordered IP meets; their
panel orders are respectively (2,4,2,2) and (2,2,2,4). Therefore they
are not thin. These failures demonstrate why neither endpoint generation
nor flag transitivity is a substitute for IP.

\subsection{Exhaustive extension of proper
residues}\label{exhaustive-extension-of-proper-residues}

The original geometry census can be replaced by a short complete search
on proper residues. The required group data are the complete TomLib
table of J1 subgroup-conjugacy classes, not a list of its geometries.

\textbf{Theorem C.Census.} Let G be finite and let \(M_{1},\ldots ,M_s\)
represent all conjugacy classes of maximal proper subgroups. For each
\(M_j\) enumerate, up to conjugation by \(M_j\), all unordered k-element
sets of involutions whose markings satisfy IP and typed Helly. Append
every involution of G outside the subgroup generated by each retained
set and apply the same tests. This finds every rank-(k+1) thin regular
residually connected geometry whose full type-preserving group is a
subgroup of G. To restrict to group G, retain exactly the extensions
generating G.

\textbf{Proof.} For any such rank-(k+1) geometry, its distinguished
marking consists of independent involutions. Omitting one gives a proper
subgroup P of its full group and hence of G, so P lies in a conjugate of
one \(M_j\). Its marked coset geometry is the corresponding rank-k
residue, thus satisfies both IP and typed Helly. Conjugation into
\(M_j\) and then within \(M_j\) puts its unordered markers in the
enumerated list. Its omitted involution lies outside P and occurs in the
extension loop. Conversely, the tested IP and Helly conditions give
exactly the asserted geometries by UC.Geometry. More formally, for each
\(M_j\) take E to be its admissible unordered k-marker sets and Z to be
pairs (T,t) with \(T\in E\), t an involution outside
\(\langle T\rangle\), and the required extension properties. Conjugation
by \(M_j\) acts on both sets, and \((T,t)\mapsto T\) is equivariant.
UC.ExtensionOrbits identifies the extension orbits over a representative
T with the stabilizer orbits in its fiber. The certificate tries the
whole fiber, so it certainly meets each such orbit; its final conjugacy
keys remove duplicate extensions. Ordering the marker set merely orders
the type labels. \(\square\)

For completeness, the residue statement used here requires no extension
of arbitrary residue automorphisms to G. The base-flag stabilizer acts
transitively and freely on the residue's chambers, and the residue is
thin and connected. Its full type-preserving group also acts freely, so
the stabilizer already is that full group. Walking inside the residue
identifies its panel generators. Thus local Helly is a necessary
condition for these genuine residues, not an arbitrary extra filter.

\textbf{Corollary C.Classification.} Up to isomorphism allowing a
permutation of types, the six geometries in Corollary C.J1 are all the
thin regular residually connected rank-four geometries with group J1.
There are none of rank at least five.

\textbf{Proof.} The exhaustive certificate executes Theorem C.Census for
\(k=3\) and \(k=4\). The retained complete table has forty subgroup
classes. Its nonzero containment entries determine seven maximal classes
of orders 660,168,120,114,110,60,42. The local counts are:

{\small % Uncaptioned data table.
\begin{longtable}[]{@{}
  >{\raggedleft\arraybackslash}p{(\linewidth - 10\tabcolsep) * \real{0.1667}}
  >{\raggedleft\arraybackslash}p{(\linewidth - 10\tabcolsep) * \real{0.1667}}
  >{\raggedleft\arraybackslash}p{(\linewidth - 10\tabcolsep) * \real{0.1667}}
  >{\raggedleft\arraybackslash}p{(\linewidth - 10\tabcolsep) * \real{0.1667}}
  >{\raggedleft\arraybackslash}p{(\linewidth - 10\tabcolsep) * \real{0.1667}}
  >{\raggedleft\arraybackslash}p{(\linewidth - 10\tabcolsep) * \real{0.1667}}@{}}
\toprule\noalign{}
\begin{minipage}[b]{\linewidth}\raggedleft
Maximal subgroup order
\end{minipage} & \begin{minipage}[b]{\linewidth}\raggedleft
Involutions
\end{minipage} & \begin{minipage}[b]{\linewidth}\raggedleft
Triple orbits
\end{minipage} & \begin{minipage}[b]{\linewidth}\raggedleft
Quadruple orbits
\end{minipage} & \begin{minipage}[b]{\linewidth}\raggedleft
Regular triple orbits
\end{minipage} & \begin{minipage}[b]{\linewidth}\raggedleft
Regular quadruple orbits
\end{minipage} \\
\midrule\noalign{}
\endhead
\bottomrule\noalign{}
\endlastfoot
660 & 55 & 60 & 588 & 42 & 6 \\
168 & 7 & 3 & 3 & 2 & 0 \\
120 & 31 & 108 & 616 & 45 & 7 \\
114 & 19 & 12 & 42 & 0 & 0 \\
110 & 11 & 2 & 4 & 0 & 0 \\
60 & 23 & 58 & 221 & 6 & 2 \\
42 & 7 & 2 & 2 & 0 & 0 \\
Total & --- & 245 & 1476 & 95 & 15 \\
\end{longtable}
}

The finite conjugation action is used here without changing its orbit
relation. Conjugation on the actual set of involutions has a permutation
image; every element of that image has an ambient conjugating lift, so
its induced subset orbits are exactly the needed conjugacy orbits.
Faithfulness of this action is not assumed. No local generated subgroup
is required to equal its containing maximal subgroup. The 95 retained
triples reduce to 70 J1-conjugacy classes. All 100016 possible
outside-involution extensions are tested; forty pass with generated
group J1, collapsing to precisely six classes. Their maximal-parabolic
order multisets are the six distinct ones in Corollary C.J1, which
already supplies six pairwise nonisomorphic constructions.

Here is the exact conjugacy test, so equality of numerical invariants is
not mistaken for isomorphism. There is one class of 1463 involutions in
J1, with centralizer order 120. Fix \(a_{0}\). For each x in an
unordered marker set T choose \(g_x\) with \(x^{g_x}=a_{0}\), enumerate
\(T^{g_xc}\) for every \(c\in C_G(a_{0})\), and take the least sorted
tuple of involution indices. These are exactly the conjugates of T
containing \(a_{0}\): every conjugator carrying x to \(a_{0}\) differs
from \(g_x\) on the right by such a c.~Consequently keys agree exactly
for inner-conjugate unordered markings. There is no assumption that all
automorphisms of J1 are inner. The six different parabolic-order
multisets prevent any further isomorphism from merging those classes,
while the complete six-class upper bound prevents any additional
geometry.

For \(k=4\), the fifteen local regular quadruple classes remain fifteen
after J1-conjugacy deduplication. All 21352 extensions by an involution
outside their generated subgroup fail the full regular-geometry tests.
This stage does not discard proper-subgroup targets. If a regular
geometry of rank \(r\ge 5\) existed for J1, fixing \(r-5\) objects of a
base chamber would yield a rank-five regular residue with group
contained in J1; Theorem C.Census would have found it. Thus all higher
ranks are excluded. This does not assert nonexistence of C-group
markings that fail typed Helly. \(\square\)

The source's successful central modifications in §8, when they generate
J1 and yield a thin regular geometry, are consequently covered by the
same exhaustive classification. The two displayed failures were checked
separately above. The construction certificate also verifies the
incidental source equality
\(N_G(\langle a,b,c\rangle )=\langle a,b,c\rangle\). Its initial vertex
figure is the universal spherical {[}3,5{]} geometry, since that
parabolic has the full Coxeter order 120; its facet is the antipodal
hemi-dodecahedral quotient of {[}5,3{]}, since the A5 quotient has order
60 and kernel the unique central order-two subgroup of
\(C_{2}\times A_{5}\). These conventional spherical identifications use
the named classical Coxeter models.
